Konduktivita závisí na \textbf{koncentraci, náboji a pohyblivosti nosičů}:
\[\boxed{\sigma=n|q|\mu}\quad[\mathrm{S/m}].\]
$n$ [m$^{-3}$] je koncentrace nosičů, $|q|$ [C] velikost jejich náboje a $\mu$ [m$^2$/(V\,s)] pohyblivost. Větší koncentrace nebo pohyblivost znamená větší konduktivitu.

Vodivostní složka ztrátového činitele je
\[\tan\delta_{\mathrm{vod}}=\frac{\sigma}{\omega\varepsilon_0\varepsilon'},\qquad\omega=2\pi f.\]
Při ostatních parametrech stejných tedy vodivostní ztráty rostou s konduktivitou a jejich ztrátový činitel klesá s frekvencí jako $1/f$.
